\begin{abstract}
A strongly lensed binary-black-hole merger whose host galaxy is identified provides a time delay measured to milliseconds, the relative and apparent amplitudes of its images, and electromagnetic constraints on the lens, and therefore a measurement of the Hubble constant $H_0$. We test whether such measurements are unbiased for a population of systems. We analyze 60 lensed binary black holes detected by the Einstein Telescope (ET) and 59 detected by an O5-like network, drawn from four galaxy-lens families, combining joint parameter estimation of both gravitational-wave (GW) images from simulated detector data with a Bayesian reconstruction of the lens from image positions, lens light and velocity dispersion. Two corrections are required. The prior on the source position must be normalized for each lens as in the population from which the systems come; otherwise the reconstruction favors steep density profiles. And the likelihood must be divided by the fraction of systems that would be detected at the assumed $H_0$, which for O5 rises steeply with $H_0$. Without them, an elliptical power-law (EPL) lens model overestimates $H_0$ by 5\% (ET) and 12\% (O5). With both, the combined estimates are $H_0/H_0^{\rm true}=1.010^{+0.008}_{-0.007}$ (ET) and $1.001\pm0.011$ (O5) with the lens models of the simulated families and $1.022\pm0.010$ and $1.010^{+0.017}_{-0.015}$ with an EPL model with external shear for every system, also when an external convergence is drawn for each system. A single system constrains $H_0$ to a median 7\% (ET) and 12\% (O5). The GW apparent distance narrows these intervals by only 3--4\%: the information comes from the delay and the lens model, so the accuracy is limited, as in time-delay cosmography with lensed quasars, by the priors on the radial slope of the lens and on the external convergence.
\end{abstract}
